\documentclass[a4,brazil]{brthesis}

% ---
% Informações de dados para CAPA e FOLHA DE ROSTO
% ---
\titulo{Um exemplo de tese escrita com \textsf{brthesis}}
\autor{Fábio Pinto Fortkamp}
\local{Florianópolis}
\data{Dezembro de 2015}
\orientador{Meu orientador}

\instituicao{%
  Universidade do Brasil
  \par
  Programa de Pós-Graduação em Engenharia Mecânica}
\tipotrabalho{Tese (Doutorado em Engenharia Mecânica)}
% O preambulo deve conter o tipo do trabalho, o objetivo, 
% o nome da instituição e a área de concentração 
\preambulo{Tese submetida ao Programa de Pós-Graduação em Engenharia Mecânica para a obtenção do título de Mestre em Engenharia Mecânica.}
% ---

% this configures PDF metadata with the above information
\configurepdf

% ---
% compila o indice
% ---
\makeindex
% ---

% ----
% Início do documento
% ----
\begin{document}

% Retira espaço extra obsoleto entre as frases.
\frenchspacing 

% ----------------------------------------------------------
% ELEMENTOS PRÉ-TEXTUAIS
% ----------------------------------------------------------
\pretextual

% ---
% Capa
% ---
\printcover
% ---

%% ---
%% Folha de rosto
%% (o * indica que haverá a ficha bibliográfica)
%% ---
%\imprimirfolhaderosto*
%% ---

% ---
% Inserir a ficha bibliografica
% ---

% Isto é um exemplo de Ficha Catalográfica, ou ``Dados internacionais de
% catalogação-na-publicação''. Você pode utilizar este modelo como referência. 
% Porém, provavelmente a biblioteca da sua universidade lhe fornecerá um PDF
% com a ficha catalográfica definitiva após a defesa do trabalho. Quando estiver
% com o documento, salve-o como PDF no diretório do seu projeto e substitua todo
% o conteúdo de implementação deste arquivo pelo comando abaixo:
%
% \begin{fichacatalografica}
%     \includepdf{fig_ficha_catalografica.pdf}
% \end{fichacatalografica}

% \include{Ficha_catalografica}
% ---

% ---
% Inserir folha de aprovação
% ---

% Isto é um exemplo de Folha de aprovação, elemento obrigatório da NBR
% 14724/2011 (seção 4.2.1.3). Você pode utilizar este modelo até a aprovação
% do trabalho. Após isso, substitua todo o conteúdo deste arquivo por uma
% imagem da página assinada pela banca com o comando abaixo:
%
% \includepdf{folhadeaprovacao_final.pdf}
%
%

%\include{Folha_de_aprovacao}
% ---

% ---
% Dedicatória
% ---
%\include{Dedicatoria}
% ---

% ---
% Agradecimentos
% ---
%\include{Agradecimentos}
% ---

% ---
% Epígrafe
% ---
%\include{Epigrafe}
% ---

% ---
% RESUMOS
% ---

% resumo em português
%\include{Resumo}

% resumo em inglês
%\include{Abstract}
% ---

% inserir lista de ilustrações
\printlof

% inserir lista de tabelas
\printlot

% inserir lista de símbolos
\printlos

% inserir o sumario
\printtoc

% ----------------------------------------------------------
% ELEMENTOS TEXTUAIS
% ----------------------------------------------------------
\textual

\chapter{Um título de capítulo}
\label{cha:chapter-title}

Vamos começar com uma equação:

\begin{equation}
  \label{eq:1}
  c^2 = a^2 + b^2
\end{equation}

\noindent que pode ser referenciada por \textsl{autoref}: \autoref{eq:1}.

\section{Um título de seção}
\label{sec:section-title}

Em português, números devem ser impressos com uma vírgula como separador decimal, como em $\num{3.14}$ ou $\SI{1.3}{\pascal}$

\subsection{Uma subseção}
\label{sec:subsection}

Veja por exemplo \autoref{fig:empty}.

\begin{figure}[!ht]
  \centering
  \includegraphics[width=12cm]{example.jpg}
  \caption{Uma figura}
  \label{fig:empty}
\end{figure}


\subsection{Referências}
\label{sec:referencias}

Podem haver referências mais explícitas \cite{bib:furlani}, ou ao longo do texto, como em \citeonline{bib:aprea}.

% ----------------------------------------------------------
% ELEMENTOS PÓS-TEXTUAIS
% ----------------------------------------------------------
\postextual

% ----------------------------------------------------------
% Referências bibliográficas
% ----------------------------------------------------------
\bibliography{refs}



\end{document}


%%% Local Variables:
%%% mode: latex
%%% TeX-master: t
%%% End:
